% ============================================================
% 05-bezier-matrix-sympy-text-code-explained.tex
% Compile with:  pdflatex -shell-escape 05-bezier-matrix-sympy-text-code-explained.tex
% Run twice for table of contents.
% ============================================================

\documentclass[12pt, a4paper]{article}

% ---- Encoding & fonts ----------------------------------------
\usepackage[T1]{fontenc}
\usepackage[utf8]{inputenc}
\usepackage{lmodern}

% ---- Page geometry -------------------------------------------
\usepackage[
  top=2.5cm, bottom=2.5cm,
  left=2.8cm, right=2.8cm
]{geometry}

% ---- Colours -------------------------------------------------
\usepackage[dvipsnames]{xcolor}
\definecolor{codebg}{RGB}{248,248,248}
\definecolor{rulecolor}{RGB}{180,180,180}
\definecolor{examplebg}{RGB}{240,248,255}
\definecolor{examplerule}{RGB}{100,149,237}
\definecolor{notebg}{RGB}{255,250,230}
\definecolor{noterule}{RGB}{210,180,50}
\definecolor{samebg}{RGB}{245,245,245}
\definecolor{samerule}{RGB}{160,160,160}

% ---- Minted --------------------------------------------------
\usepackage{minted}
\setminted[julia]{
  bgcolor    = codebg,
  frame      = leftline,
  framerule  = 1.5pt,
  rulecolor  = rulecolor,
  linenos    = false,
  fontsize   = \small,
  breaklines = true,
  tabsize    = 4,
  style      = friendly,
}
\setminted[text]{
  bgcolor    = codebg,
  fontsize   = \small,
  breaklines = true,
}

% ---- Math ----------------------------------------------------
\usepackage{amsmath}
\usepackage{amssymb}

% ---- TikZ ----------------------------------------------------
\usepackage{tikz}
\usetikzlibrary{shapes.geometric, arrows.meta, positioning, backgrounds}

% ---- Boxes ---------------------------------------------------
\usepackage{mdframed}

\newmdenv[
  backgroundcolor  = examplebg,
  linecolor        = examplerule, linewidth=1.5pt,
  leftline=true, rightline=false, topline=false, bottomline=false,
  innerleftmargin=10pt, innerrightmargin=8pt,
  innertopmargin=6pt, innerbottommargin=6pt,
  skipabove=6pt, skipbelow=4pt,
]{examplebox}

\newmdenv[
  backgroundcolor  = notebg,
  linecolor        = noterule, linewidth=1.5pt,
  leftline=true, rightline=false, topline=false, bottomline=false,
  innerleftmargin=10pt, innerrightmargin=8pt,
  innertopmargin=6pt, innerbottommargin=6pt,
  skipabove=6pt, skipbelow=4pt,
]{notebox}

\newmdenv[
  backgroundcolor  = samebg,
  linecolor        = samerule, linewidth=1.5pt,
  leftline=true, rightline=false, topline=false, bottomline=false,
  innerleftmargin=10pt, innerrightmargin=8pt,
  innertopmargin=6pt, innerbottommargin=6pt,
  skipabove=6pt, skipbelow=4pt,
]{samebox}

% ---- Hyperlinks ----------------------------------------------
\usepackage[colorlinks=true, linkcolor=black, urlcolor=MidnightBlue]{hyperref}

% ---- Misc ----------------------------------------------------
\usepackage{booktabs}
\usepackage{needspace}
\usepackage{parskip}
\usepackage{titlesec}
\usepackage{fancyhdr}
\usepackage{datetime2}

% ---- Section formatting --------------------------------------
\titleformat{\section}
  {\large\bfseries\color{MidnightBlue}}
  {\thesection.}{0.6em}{}[\vspace{-4pt}\rule{\linewidth}{0.4pt}]
\titleformat{\subsection}
  {\normalsize\bfseries\color{BrickRed}}
  {\thesubsection.}{0.5em}{}
\titleformat{\subsubsection}
  {\normalsize\bfseries}
  {\thesubsubsection.}{0.5em}{}

% ---- Header / footer -----------------------------------------
\pagestyle{fancy}
\fancyhf{}
\rhead{\textcolor{gray}{\small 05-bezier-matrix-sympy-text-code-explained}}
\lhead{\textcolor{gray}{\small B\'ezier Matrix Symbolic --- Code Explanation}}
\cfoot{\thepage}
\renewcommand{\headrulewidth}{0.3pt}

% ---- Helpers -------------------------------------------------
\newcommand{\cd}[1]{\texttt{#1}}

% ---- Title ---------------------------------------------------
\title{%
  \vspace{1cm}
  {\LARGE\bfseries B\'ezier Curve --- Matrix Method (Symbolic)}\\[0.4em]
  {\large $P(u) = U \cdot M \cdot G$ \quad via Symbolics.jl}\\[0.4em]
  {\large Code Explanation for Julia Beginners}\\[0.4em]
  {\normalsize\texttt{05-bezier-matrix-sympy-text.jl}}
}
\author{E.R.\ Nurwijayadi}
\date{\DTMtoday}


% =============================================================
\begin{document}
% =============================================================

\maketitle
\thispagestyle{empty}

\begin{center}
\textit{Directed and curated by E.R.\ Nurwijayadi.
Code and explanations AI-generated.}
\end{center}

\begin{center}
\begin{minipage}{0.85\textwidth}
\small
This document explains the Julia code in
\cd{05-bezier-matrix-sympy-text.jl} for readers new to Julia.
Script~05 is a deliberate hybrid: it takes the straight
output structure of script~04 and substitutes exact
\cd{Rational\{Int\}} arithmetic throughout, using
Symbolics.jl only for polynomial formatting.
This document focuses on what changed relative to
scripts~02 and~04, which are its two parents.
Shared components are listed by reference.
\end{minipage}
\end{center}

\vspace{1em}\hrule\vspace{1.5em}
\tableofcontents
\newpage


% =============================================================
\needspace{8\baselineskip}
\section{Type Map: How the Pieces Fit Together}
% =============================================================

\textit{Script~05 combines the exact arithmetic engine of
script~02 with the straight print structure of script~04.
In the diagram, grey boxes are inherited unchanged,
white boxes are new or modified.}

\bigskip

\tikzset{
  structbox/.style={
    rectangle, rounded corners=4pt,
    draw=MidnightBlue, line width=1.5pt,
    fill=examplebg,
    text width=3.6cm, align=center,
    font=\small\bfseries, inner sep=8pt
  },
  funcbox/.style={
    rectangle, rounded corners=2pt,
    draw=gray!60, line width=1pt,
    fill=white,
    text width=3.6cm, align=center,
    font=\small, inner sep=7pt
  },
  samefunc/.style={
    rectangle, rounded corners=2pt,
    draw=samerule, line width=1pt,
    fill=samebg,
    text width=3.6cm, align=center,
    font=\small, inner sep=7pt
  },
  casbox/.style={
    rectangle, rounded corners=4pt,
    draw=BrickRed!60, line width=1.5pt,
    fill=notebg,
    text width=3.6cm, align=center,
    font=\small\bfseries, inner sep=8pt
  },
  entrybox/.style={
    rectangle, rounded corners=4pt,
    draw=BrickRed, line width=1.5pt,
    fill=notebg,
    text width=2.4cm, align=center,
    font=\small\bfseries, inner sep=8pt
  },
  myarrow/.style={-Stealth, thick, draw=gray!70},
  bluearrow/.style={-Stealth, thick, draw=MidnightBlue},
  redarrow/.style={-Stealth, thick, draw=BrickRed!80},
  arrowlabel/.style={font=\footnotesize\itshape, fill=white, inner sep=2pt}
}

\begin{center}
\begin{tikzpicture}[node distance=1.4cm and 2.8cm]

  % -- Symbolics.jl (external) ------------------------------
  \node[casbox] (symjl) {%
    \texttt{Symbolics.jl}\\[3pt]
    \footnotesize\normalfont
    \texttt{@variables u}\\
    \texttt{expand()}, \texttt{string()}\\
    \textit{for fmt\_poly only}
  };

  % -- BezierMatrix struct (modified) -----------------------
  \node[structbox, right=3.2cm of symjl] (bm) {%
    \texttt{BezierMatrix}\\[3pt]
    \footnotesize\normalfont
    \texttt{G, M, MG}\\
    \texttt{::Matrix\{Rational\{Int\}\}}\\
    \texttt{axes, n, degree, dim}
  };

  % -- derive_M + evaluate_all (modified type) --------------
  \node[funcbox, below left=1.6cm and 0.3cm of bm] (core) {%
    \texttt{derive\_M()} \textbf{[mod]}\\
    \texttt{evaluate\_all()} \textbf{[mod]}\\[2pt]
    \footnotesize\normalfont
    \textit{same logic, type changed}\\
    \textit{to \texttt{Rational\{Int\}}}
  };

  % -- fmt_rat + fmt_poly (new/from script 02) --------------
  \node[funcbox, below right=1.6cm and 0.3cm of bm] (fmt) {%
    \texttt{fmt\_rat()} \textbf{[from 02]}\\
    \texttt{fmt\_poly()} \textbf{[new]}\\[2pt]
    \footnotesize\normalfont
    \textit{rational formatter}\\
    \textit{symbolic poly builder}
  };

  % -- table printers (modified type) -----------------------
  \node[samefunc, below=1.3cm of core] (tbl) {%
    \textbf{Table printers}\\[2pt]
    \footnotesize\normalfont
    \texttt{print\_table} (unchanged)\\
    \texttt{print\_labeled\_matrix}\\
    \textit{(type: Rational\{Int\})}
  };

  % -- section() (from script 04) ---------------------------
  \node[samefunc, below=1.3cm of fmt] (sec) {%
    \texttt{section()} \textbf{[from 04]}\\[2pt]
    \footnotesize\normalfont
    \textit{unchanged from script~04}
  };

  % -- print functions (from script 04, minor mods) ---------
  \node[funcbox, below=1.3cm of tbl, xshift=2.2cm] (print) {%
    \textbf{Print functions}\\[2pt]
    \footnotesize\normalfont
    \texttt{print\_G/M/MG} (type mod)\\
    \texttt{print\_polynomials} (mod)\\
    \texttt{print\_eval\_table} (mod)\\
    \texttt{print\_all} (unchanged)
  };

  % -- main (unchanged) -------------------------------------
  \node[entrybox, right=2.0cm of bm] (main) {%
    \texttt{main()}\\[2pt]
    \footnotesize\normalfont\normalcolor
    \textit{entry point}\\
    \textit{unchanged}
  };

  % -- Arrows -----------------------------------------------
  \draw[myarrow]   (symjl)  -- node[arrowlabel, above]{used by} (fmt);
  \draw[bluearrow] (core)   -- node[arrowlabel, left]{builds} (bm);
  \draw[myarrow]   (bm)     -- node[arrowlabel, above left]{calls} (core);
  \draw[myarrow]   (print)  -- node[arrowlabel, left]{calls}  (tbl);
  \draw[myarrow]   (print)  -- node[arrowlabel, right]{calls} (sec);
  \draw[myarrow]   (print)  -- node[arrowlabel, above]{calls} (fmt);
  \draw[myarrow]   (print)  -- node[arrowlabel, above]{reads} (bm);
  \draw[myarrow]   (print)  |- node[arrowlabel, right,pos=0.8]{calls} (core);
  \draw[redarrow]  (main)   -- node[arrowlabel, above]{calls} (bm);
  \draw[redarrow]  (main)   |- node[arrowlabel, right,pos=0.75]{calls} (print);

\end{tikzpicture}
\end{center}

\bigskip
\begin{notebox}
\textbf{Script~05 is a directed combination of two parents.}

From \textbf{script~04}: the straight print structure ---
\cd{section()}, five focused print functions, \cd{print\_all}.

From \textbf{script~02}: exact rational arithmetic ---
\cd{Rational\{Int\}} types, \cd{fmt\_rat}, Symbolics.jl for
polynomial display.

The only genuinely new code is \cd{fmt\_poly}, which bridges
the two parents: it takes a \cd{Rational\{Int\}} column from
\cd{MG} and turns it into a Symbolics.jl symbolic expression.
\end{notebox}


% =============================================================
\needspace{8\baselineskip}
\section{What Is Shared (by Reference)}
% =============================================================

\begin{samebox}
\textbf{From script~04 --- identical or structurally identical:}

\medskip
\begin{tabular}{lp{8.5cm}}
\toprule
Component & Notes \\
\midrule
\cd{section()} & Identical. \\
\cd{print\_table} & Identical. \\
\cd{print\_all} signature & Takes \cd{Vector\{Rational\{Int\}\}}
  instead of \cd{Vector\{Float64\}} --- otherwise identical. \\
\cd{ruler} & Identical. \\
\cd{main} & Identical logic. \\
\bottomrule
\end{tabular}

\medskip
\textbf{From script~02 --- same concept, same name:}

\medskip
\begin{tabular}{lp{8.5cm}}
\toprule
Component & Notes \\
\midrule
\cd{@variables u} & Identical declaration. \\
\cd{Rational\{Int\}} control points & Same typed matrix literal. \\
\cd{Rational\{Int\}} u\_values & Same exact fraction list. \\
\cd{fmt\_rat} & Same two-method dispatch function
  (named \cd{fmt\_rational} in script~02). \\
\bottomrule
\end{tabular}
\end{samebox}


% =============================================================
\needspace{8\baselineskip}
\section{The \texttt{BezierMatrix} Struct --- Type Change}
% =============================================================

\begin{minted}{julia}
struct BezierMatrix
    G      ::Matrix{Rational{Int}}
    axes   ::Vector{String}
    n      ::Int
    degree ::Int
    dim    ::Int
    M      ::Matrix{Rational{Int}}
    MG     ::Matrix{Rational{Int}}
end
\end{minted}

The struct is identical to script~04 except that the three
matrix fields changed from \cd{Matrix\{Float64\}} to
\cd{Matrix\{Rational\{Int\}\}}.

This single type change propagates everywhere:
all arithmetic on these matrices now uses exact fraction
arithmetic automatically, because Julia dispatches
\cd{+}, \cd{*}, etc.\ based on the element type.

\begin{notebox}
\textbf{Why does changing the type of \cd{G} matter?}

The constructor computes \cd{MG = M * points}.
If \cd{points} is \cd{Matrix\{Float64\}}, the result is
\cd{Matrix\{Float64\}} --- floating-point multiplication.
If \cd{points} is \cd{Matrix\{Rational\{Int\}\}}, the result
is \cd{Matrix\{Rational\{Int\}\}} --- exact rational multiplication.
Julia infers the output type from the input types automatically.
No code change in the constructor is needed beyond the type
annotation on the struct field.
\end{notebox}


% =============================================================
\needspace{8\baselineskip}
\section{\texttt{derive\_M} and \texttt{evaluate\_all} --- Type Change Only}
% =============================================================

\needspace{5\baselineskip}
\subsection{\texttt{derive\_M}}

\begin{minted}{julia}
function derive_M(deg::Int)::Matrix{Rational{Int}}
    n = deg
    M = zeros(Rational{Int}, n+1, n+1)
    for i in 0:n
        m = n - i
        for j in 0:m
            k = i + j
            M[n-k+1, i+1] = binomial(n,i) * binomial(m,j) * (-1)^j
        end
    end
    return M
end
\end{minted}

Two changes from script~04:

\begin{enumerate}
  \item Return type annotation: \cd{::Matrix\{Rational\{Int\}\}}
        instead of \cd{::Matrix\{Float64\}}.
  \item Initialisation: \cd{zeros(Rational\{Int\}, n+1, n+1)}
        instead of \cd{zeros(Float64, n+1, n+1)}.
\end{enumerate}

\cd{zeros(Rational\{Int\}, n+1, n+1)} creates a matrix
filled with \cd{0//1} (the rational zero).
The loop body is identical: \cd{binomial(n,i)},
\cd{binomial(m,j)}, and \cd{(-1)\^{}j} all return
plain integers, and assigning an integer to a
\cd{Rational\{Int\}} slot auto-promotes it to
\cd{binomial\_result // 1}.

\begin{examplebox}
\textbf{For $n=4$, $i=0$, $j=0$:}

\cd{binomial(4,0) * binomial(4,0) * (-1)\^{}0 = 1 * 1 * 1 = 1}

Assigned to \cd{M[5, 1]} (row \cd{n-k+1 = 5}, col \cd{i+1 = 1}).

Stored as \cd{1//1} in \cd{Matrix\{Rational\{Int\}\}}.

Displayed as \cd{1} by \cd{fmt\_rat} (denominator is 1).
\end{examplebox}

\needspace{5\baselineskip}
\subsection{\texttt{evaluate\_all}}

\begin{minted}{julia}
function evaluate_all(bm::BezierMatrix,
    u_values::Vector{Rational{Int}})::Matrix{Rational{Int}}
    n    = bm.degree
    Uall = [u_val^(n-k) for u_val in u_values, k in 0:n]
    Uall * bm.MG
end
\end{minted}

Three type changes from script~04:
\begin{itemize}
  \item Argument: \cd{u\_values::Vector\{Rational\{Int\}\}}
  \item Return annotation: \cd{::Matrix\{Rational\{Int\}\}}
  \item The local variable \cd{u} from script~04 is renamed
        \cd{u\_val} to avoid shadowing the Symbolics.jl
        symbolic variable \cd{u} declared at the top of
        the file with \cd{@variables u}.
\end{itemize}

\begin{notebox}
\textbf{Why rename \cd{u} to \cd{u\_val}?}

At the top of the file, \cd{@variables u} creates a
symbolic object named \cd{u}.
If the loop variable in the comprehension were also named
\cd{u}, it would \emph{shadow} the symbolic \cd{u} inside
the comprehension body --- meaning \cd{u\^{}(n-k)} would
use the rational loop value, not the symbolic variable.
In this function that would actually be correct,
but it would be confusing and fragile.
Renaming to \cd{u\_val} keeps the two \cd{u}s
unambiguously separate throughout the file.
\end{notebox}


% =============================================================
\needspace{8\baselineskip}
\section{\texttt{fmt\_rat} --- The Rational Formatter}
% =============================================================

\begin{minted}{julia}
fmt_rat(r::Rational{Int}) =
    isone(denominator(r)) ? string(numerator(r)) : string(r)
fmt_rat(n::Integer) = string(n)
\end{minted}

This is identical in purpose and structure to
\cd{fmt\_rational} from script~02, renamed for brevity.

The two-method multiple dispatch covers:
\begin{itemize}
  \item A \cd{Rational\{Int\}} value: print as \cd{"5"} if
        denominator is 1, else as \cd{"269//256"}.
  \item A plain \cd{Integer}: just convert to string.
\end{itemize}

It is used in three places:
\begin{enumerate}
  \item \cd{print\_labeled\_matrix} --- to format every cell
        of $G$, $M$, and $M \cdot G$.
  \item \cd{print\_eval\_table} --- to format the $u$ value
        column and the exact result columns.
\end{enumerate}


% =============================================================
\needspace{8\baselineskip}
\section{\texttt{fmt\_poly} --- Building a Symbolic Expression}
% =============================================================

\textit{This is the only genuinely new function in the script.
It is short but dense --- four lines that do something
quite different from any formatting function in the
previous scripts.}

\begin{minted}{julia}
function fmt_poly(mg_col::Vector{Rational{Int}}, deg::Int)::String
    expr = sum(
        mg_col[k+1] * u^(deg-k)
        for k in 0:deg
        if mg_col[k+1] != 0
    ; init = 0*u^0)
    string(expand(expr))
end
\end{minted}

\needspace{5\baselineskip}
\subsection{What it does}

Instead of building a string directly from the coefficient
vector (as script~04's \cd{fmt\_poly} did), this function:
\begin{enumerate}
  \item Builds a \textbf{Symbolics.jl symbolic expression}
        by summing terms of the form
        \cd{coefficient * u\^{}power}.
  \item Calls \cd{expand()} to canonicalise the expression.
  \item Calls \cd{string()} to convert the Symbolics
        expression to its printed representation.
\end{enumerate}

\needspace{5\baselineskip}
\subsection{The generator sum with \texttt{if} filter}

\begin{minted}{julia}
sum(
    mg_col[k+1] * u^(deg-k)
    for k in 0:deg
    if mg_col[k+1] != 0
; init = 0*u^0)
\end{minted}

This is \cd{sum} applied to a \textbf{generator expression}
with an inline \cd{if} filter.
It is equivalent to:

\begin{minted}{julia}
total = 0*u^0          # initial value
for k in 0:deg
    if mg_col[k+1] != 0
        total += mg_col[k+1] * u^(deg-k)
    end
end
\end{minted}

The generator computes \cd{mg\_col[k+1] * u\^{}(deg-k)}
for each \cd{k} where the coefficient is non-zero,
and \cd{sum} accumulates them.

\begin{examplebox}
\textbf{For the $x$ column: \cd{mg\_col = [13//1, -32//1, 24//1, 0//1, 0//1]},
\cd{deg = 4}:}

\medskip
\begin{tabular}{lllll}
\toprule
\cd{k} & \cd{mg\_col[k+1]} & non-zero? & term built \\
\midrule
0 & \cd{13//1}  & yes & \cd{13//1 * u\^{}4} \\
1 & \cd{-32//1} & yes & \cd{-32//1 * u\^{}3} \\
2 & \cd{24//1}  & yes & \cd{24//1 * u\^{}2} \\
3 & \cd{0//1}   & no  & skipped \\
4 & \cd{0//1}   & no  & skipped \\
\bottomrule
\end{tabular}

\medskip
\cd{expr} = symbolic expression for
$\frac{13}{1}u^4 + \frac{-32}{1}u^3 + \frac{24}{1}u^2$

After \cd{expand()}: same expression, collected and ordered.

\cd{string(expand(expr))}:
\cd{"(24//1)*(u\^{}2) - (32//1)*(u\^{}3) + (13//1)*(u\^{}4)"}

Note the term ordering: Symbolics.jl orders from lowest to
highest power, which is why $u^2$ appears before $u^4$.
\end{examplebox}

\needspace{5\baselineskip}
\subsection{The \texttt{init} keyword argument}

\begin{minted}{julia}
sum(...; init = 0*u^0)
\end{minted}

\cd{init} provides the \textbf{starting value} for the
accumulation.
Without it, if \emph{all} coefficients are zero (e.g.\ for
a zero polynomial), \cd{sum} would have nothing to add and
would throw an error.

\cd{0*u\^{}0} creates the Symbolics expression $0 \cdot u^0$,
which simplifies to $0$ --- a symbolic zero.
Using a symbolic zero rather than a plain \cd{0} ensures
the result of \cd{sum} is always a Symbolics expression,
even when all terms are filtered out.

\begin{examplebox}
\textbf{For the $z$ column: \cd{mg\_col = [0//1, 0//1, 0//1, 0//1, 1//1]}:}

\medskip
Only \cd{k=4} passes the filter:
\cd{1//1 * u\^{}0 = 1//1 * 1 = 1//1}

\cd{sum} accumulates: \cd{0*u\^{}0 + 1//1 = 1//1}

\cd{expand(1//1)} $\to$ symbolic \cd{1}

\cd{string(1)} $\to$ \cd{"1"}

Printed as: \cd{z(u)  =  1}
\end{examplebox}

\begin{notebox}
\textbf{Why \cd{0*u\^{}0} and not just \cd{0}?}

If \cd{init = 0} (a plain integer) and the first term
added is \cd{13//1 * u\^{}4} (a Symbolics expression),
Julia would try to compute \cd{0 + (Symbolics expression)}.
This works in Julia because Symbolics.jl defines addition
with integers, but using \cd{0*u\^{}0} is more explicit:
it signals that the accumulator should be a Symbolics
expression from the start, avoiding any type ambiguity.
\end{notebox}


% =============================================================
\needspace{8\baselineskip}
\section{\texttt{print\_labeled\_matrix} --- Type Change}
% =============================================================

\begin{minted}{julia}
function print_labeled_matrix(
        mat        ::Matrix{Rational{Int}},
        row_labels ::Vector{String},
        col_labels ::Vector{String})
    rows = [[fmt_rat(mat[r,c]) for c in 1:size(mat,2)]
            for r in 1:size(mat,1)]
    print_table(rows, col_labels; row_labels=row_labels)
end
\end{minted}

Two changes from script~04:
\begin{itemize}
  \item \cd{mat} type: \cd{Matrix\{Rational\{Int\}\}} instead
        of \cd{Matrix\{Float64\}}.
  \item Formatting function: \cd{fmt\_rat} instead of
        \cd{fmt\_fn} (the configurable keyword argument
        is gone, replaced by the single fixed formatter
        appropriate for rational values).
\end{itemize}

Removing the \cd{fmt\_fn} keyword argument is possible here
because \cd{Rational\{Int\}} values always use \cd{fmt\_rat}
--- there is no longer a reason to offer an alternative formatter.


% =============================================================
\needspace{8\baselineskip}
\section{The Print Functions --- What Changed}
% =============================================================

\needspace{5\baselineskip}
\subsection{\texttt{print\_G}, \texttt{print\_M}, \texttt{print\_MG}}

These three functions are structurally identical to
script~04.
The only difference is that \cd{print\_labeled\_matrix}
now accepts \cd{Matrix\{Rational\{Int\}\}}, so the call
sites work without any change to the function bodies.

\begin{minted}{julia}
function print_G(bm::BezierMatrix)
    section("G  - Geometry Matrix (Control Points)")
    println()
    print_labeled_matrix(bm.G, ["P$(i-1)" for i in 1:bm.n], bm.axes)
end
\end{minted}

\textbf{Output:} identical to script~04 because all
values in $G$ are whole numbers and \cd{fmt\_rat} prints
\cd{4//1} as \cd{4}.

\needspace{5\baselineskip}
\subsection{\texttt{print\_polynomials} --- uses \texttt{fmt\_poly}}

\begin{minted}{julia}
function print_polynomials(bm::BezierMatrix)
    section("Resulting Polynomials  (via Symbolics.jl)")
    println()
    for (d, axis) in enumerate(bm.axes)
        println("  $(axis)(u)  =  $(fmt_poly(bm.MG[:,d], bm.degree))")
    end
end
\end{minted}

The loop body is identical to script~04 but calls
\cd{fmt\_poly} (the new Symbolics-based formatter) instead
of script~04's hand-written \cd{fmt\_poly}.
\cd{bm.MG[:,d]} passes a \cd{Vector\{Rational\{Int\}\}}
column to the new \cd{fmt\_poly}, which builds the symbolic
expression and returns its string representation.

\textbf{Output:}
\begin{minted}{text}
  Resulting Polynomials  (via Symbolics.jl)
  ----------------------------------------------------------------

  x(u)  =  (24//1)*(u^2) - (32//1)*(u^3) + (13//1)*(u^4)
  y(u)  =  (16//1)*u - (48//1)*(u^2) + (64//1)*(u^3) - (28//1)*(u^4)
  z(u)  =  1
\end{minted}

\needspace{5\baselineskip}
\subsection{\texttt{print\_eval\_table} --- rational arithmetic path}

\begin{minted}{julia}
function print_eval_table(bm::BezierMatrix,
                          u_values::Vector{Rational{Int}})
    section("Evaluation Table  -  P(u) = U * (M * G)")
    pts  = evaluate_all(bm, u_values)
    rows = Vector{Vector{String}}()
    for (r, u_val) in enumerate(u_values)
        row = [fmt_rat(u_val)]
        for c in 1:bm.dim
            exact = pts[r, c]
            dec   = @sprintf("%.4f", Float64(exact))
            push!(row, "$(fmt_rat(exact))  ($dec)")
        end
        push!(rows, row)
    end
    println()
    print_table(rows, vcat(["u"], ["$(ax)(u)" for ax in bm.axes]);
                style=:outline)
end
\end{minted}

Compare to script~04's equivalent:

\begin{center}
\begin{tabular}{lp{4.5cm}p{4.5cm}}
\toprule
Step & Script~04 & Script~05 \\
\midrule
$u$ value display &
  \cd{fmt(u)} --- float to string &
  \cd{fmt\_rat(u\_val)} --- rational to string \\[4pt]
Results from \cd{evaluate\_all} &
  \cd{Matrix\{Float64\}} &
  \cd{Matrix\{Rational\{Int\}\}} \\[4pt]
Exact value per cell &
  \cd{fmt\_frac(v)} --- \cd{rationalize()} to recover fraction &
  \cd{fmt\_rat(exact)} --- fraction is already exact \\[4pt]
Decimal conversion &
  \cd{@sprintf("\%.4f", v)} --- float already available &
  \cd{@sprintf("\%.4f", Float64(exact))} --- convert rational to float \\
\bottomrule
\end{tabular}
\end{center}

The key difference is in the exact value column.
Script~04 computed with floats and used \cd{rationalize()}
to \emph{recover} the fraction.
Script~05 never loses the fraction --- \cd{exact} is already
a \cd{Rational\{Int\}}, so \cd{fmt\_rat} just formats it directly.

\begin{examplebox}
\textbf{Tracing one cell: $x$ at $u = \tfrac{1}{2}$}

\cd{u\_val = 1//2} (exact rational, displayed as \cd{"1//2"})

\cd{pts[3, 1] = 45//16} (from \cd{evaluate\_all}, exact rational)

\cd{fmt\_rat(45//16)} $\to$ \cd{"45//16"}

\cd{Float64(45//16) = 2.8125}

\cd{@sprintf("\%.4f", 2.8125)} $\to$ \cd{"2.8125"}

Cell string: \cd{"45//16  (2.8125)"}
\end{examplebox}


% =============================================================
\needspace{8\baselineskip}
\section{Summary: The Hybrid Design}
% =============================================================

Script~05 is the shortest in terms of new code ---
but it demonstrates the most important design principle
of the series: \textbf{composability}.

The matrix method (scripts~03 and~04) and exact rational
arithmetic (script~02) are independent choices.
Script~05 shows that combining them requires changing
only the type annotations and the one formatting function
that bridges the two systems (\cd{fmt\_poly}).
Everything else --- the matrix derivation logic, the
evaluation batch, the table printer, the print structure
--- works unchanged.

\begin{center}
\begin{tabular}{lll}
\toprule
Script & Method & Arithmetic \\
\midrule
01 & Bernstein algebraic & \cd{Float64} \\
02 & Bernstein algebraic & \cd{Rational\{Int\}} + Symbolics.jl \\
03 & Matrix (pedagogical) & \cd{Float64} \\
04 & Matrix (straight) & \cd{Float64} \\
05 & Matrix (straight) & \cd{Rational\{Int\}} + Symbolics.jl \\
\bottomrule
\end{tabular}
\end{center}

\vspace{2em}\hrule\vspace{0.5em}
{\small\textcolor{gray}{%
Code explanation for \texttt{05-bezier-matrix-sympy-text.jl}
\textbar{} Directed and curated by E.R.\ Nurwijayadi
\textbar{} Code and explanations AI-generated}}

\end{document}
